\subsection{Sprint 3}

Durante a primeira semana da Sprint, realizamos uma sessão de \textit{Planning Poker}, que é uma estratégia para estimar o esforço necessário no desenvolvimento das funcionalidades do aplicativo. Desde o início da Sprint 3, eu já sabia qual seria minha tarefa: finalizar a coleta de tempo de tela. Após a sessão de planejamento, comecei a pesquisar artigos sobre assincronicidade para me auxiliar nessa tarefa. Na Sprint 2, identifiquei um problema em que o usuário não tinha um tempo de notificação definido ao criar sua conta, e devido aos requisitos do \texttt{Pausable\_Timer}, não era possível passar um valor nulo. Decidi pedir ajuda ao \ac{ages} IV Vitor Hugo Garcez Barcellos dos Santos, pois ele havia estudado bastante sobre o assunto em sua última passagem pela \ac{ages}.

Marcamos um horário para que eu pudesse mostrar a ele o código e os erros que estavam ocorrendo. Logo de início, percebemos que, com alguns métodos implementados no final da Sprint 2, o problema de assincronicidade se tornou trivial, e seria necessário apenas refatorar partes do código. Além disso, constatamos que, com essas alterações, não seria necessário salvar o tempo dos \textit{timers} no banco, pois o código poderia ser executado em segundo plano e, quando o tempo selecionado pelo usuário expirasse, poderíamos chamar uma função.

Após essa conversa com Vitor, verificamos que o \textit{timer} da tela inicial não era retomado após o usuário desbloquear a tela. Ao analisar a documentação do \textit{plugin}, constatamos que várias pessoas também enfrentavam esse mesmo problema. Decidimos deixar essa funcionalidade como uma pendência para a última entrega, pois a tarefa de coleta de tempo de tela já estava concluída. Vitor me pediu para fazer a integração do código que envia a notificação, busca uma frase e cria um mecanismo de agendamento da notificação.

Ao analisar essas partes do código, deparei-me com as seguintes situações: o código de notificações não havia sido testado em nenhum ambiente Android e o mecanismo de agendamento de notificação ainda não havia sido implementado. Iniciei o desenvolvimento do mecanismo de agendamento, pois era o problema menos complexo e pude resolvê-lo rapidamente. Em seguida, solicitei ajuda ao \ac{ages} I Luis Felipe Copetti Rossi, que havia criado o código para chamar a notificação, para compreender melhor o código e resolver as pendências relacionadas ao Android. Após analisar a documentação da biblioteca, identificamos que seria necessário fazer algumas alterações no Gradle. Assim que adicionamos as dependências solicitadas na documentação, a integração foi concluída.

Durante a fase de \textit{Freezing}, acabei desenvolvendo um \textit{hotfix} para uma requisição no aplicativo que envia notificações. Também acompanhei o restante dos \textit{\acp{pr}} e como estavam resolvendo os conflitos e pude aprender bastante.